\begin{center}\bfseries 22.\end{center}
\noindent\hspace*{\parindent}Bei den Versuchen, auch für andere Zahlen den biquadratischen Charakter zu ermitteln, sah sich \textsc{Gauss} zwar durch Induktion zu mancherlei Einzelresultaten geführt, beim Beweise derselben wie in der Erkenntnis allgemeinerer Sätze aber durch Schwierigkeiten gehemmt, die er in keiner Weise mit
\newpage\SourceLeaf{62}\ArithmeticHead{55}{R}{KUBISCHE UND BIQUADRATISCHE RESTE. KOMPLEXE GANZE ZAHLEN.}
\noindent den Mitteln der rationalen Zahlentheorie zu überwinden vermochte. Da kam ihm eine erlösende Eingebung, die zugleich die Grundlage des wichtigsten Gebietes der ganzen neueren Zahlentheorie, der Arithmetik der Zahlkörper, geworden ist: »mox vero comperimus, principia Arithmeticae hactenus usitata ad theoriam generalem neutiquam sufficere, quin potius hanc necessario postulare, ut campus Arithmeticae sublimioris infinities quasi promoveatur« (W. II, S. 67). Er erfand die Theorie der komplexen ganzen Zahlen.

Über die Berechtigung komplexer Grössen hatte \textsc{Gauss} schon zur Zeit der D. A. sich eigene Gedanken gebildet, die er im art. 91 derselben gelegentlich zu entwickeln in Aussicht stellt — ein im Jahre 1831 in der Anzeige der \textit{Comm. secunda} (W. II, S. 174—178) eingelöstes Versprechen — und die Lehre von der Kreisteilung, mehr noch seine Untersuchungen über die lemniskatischen Funktionen hatten sie ihm immer zwingender aufgedrängt; vielleicht, dass in der Notiz Nr. 95 T. vom Oktober 1798 unter dem novus in analysi campus die Theorie der Funktionen einer komplexen Variabeln verstanden werden muss.

So lag es ihm nahe, auch komplexe ganze Zahlen in Betracht zu ziehen. Man geht aber wohl nicht fehl, wenn man für den Gedanken, sie zur Grundlage der kubischen und biquadratischen Reste zu machen, die Quelle in der Beziehung erblickt, welche zwischen dem biquadratischen Charakter der Zahl 2 und der Darstellung \(p=a^2+b^2\), ebenso zwischen ihrem kubischen Charakter und der Darstellung \(4p=a^2+27b^2\) (Nr. 133 T.) d. h. den Darstellungen von \(p\) als Norm komplexer ganzer Zahlen von \textsc{Gauss} erkannt worden war. Die Verwendung dieser komplexen Faktoren von \(p\) tritt grundlegend auch in den handschriftlichen Fragmenten über kubische und biquadratische Reste zu Tage, welche W. VIII, S. 5—14 abgedruckt und von \textsc{Fricke} erläutert sind und etwa aus derselben Zeit stammen, wie die Tagebuchnotizen Nr. 130—133, d. h. aus dem Februar 1807.

In diese Zeit fallen dann auch \textsc{Gauss}’ Versuche, zu neuen Beweisen des quadratischen Fundamentaltheorems zu gelangen, in der Hoffnung, dass diese auch auf das höhere Gebiet ausdehnbar sein möchten, eine Hoffnung, die ihn nicht trog, denn sie führten ihn zum \textsc{Gauss}schen Lemma, dessen Erweiterung auf das Gebiet der komplexen Zahlen \(a+bi\) die Grundlage seiner bezüglichen Beweisführungen bildet.
\newpage\SourceLeaf{63}\ArithmeticHead{56}{L}{BACHMANN, ÜBER GAUSS’ ZAHLENTHEORETISCHE ARBEITEN.}
\noindent\hspace*{\parindent}Neben induktiven Studien beschäftigten nun \textsc{Gauss} mannigfache Versuche, den sogenannten Dezidenten einer Zahl für einen gegebenen Modul zu berechnen bezw. das gesamte Restsystem eines Moduls in geeignete Viertel einzuteilen und dgl. mehr, wie die von \textsc{Schering} erläuterten und in die Zeit nach 1811 datierten Fragmente (W. II, S. 313—85) und andere, ihnen zeitlich wohl voraufgehende handschriftliche Notizen des Nachlasses (in Ec 3, abgedruckt W. X1, S. 56, und im Handbuch 18, Bd, vom Oktober 1805, S. 248—256) beweisen. Wohl gelegentlich solcher Versuche wurde \textsc{Gauss} auch zu dem Satze geführt, der W. X1, S. 51 aus seinem Nachlass abgedruckt ist, und der eng mit einem andern von M. A. \textsc{Stern}\footnote{Göttinger Nachrichten 1869, S. 330. Die \textsc{Stern}sche Notiz wird weiter unten als Anhang zu dem vorliegenden Aufsatze abgedruckt.}\ mitgeteilten \textsc{Gauss}schen Satze zusammenhängt.

Es muss Wunder nehmen, dass \textsc{Gauss} nirgends in seinem \textit{Tagebuche} der Einführung der komplexen Zahlen, dieses von ihm selbst für so wichtig erkannten Fortschrittes der Arithmetik, Erwähnung tut. Man könnte geneigt sein, die Nummern 144, 145 T. vom 23. Oktober 1813 dahin zu deuten; doch könnte damit unmöglich der erste Gebrauch komplexer Zahlen gemeint sein, der, wie bemerkt, viel früher datiert, vielmehr nur etwa ihre systematische Grundlegung für die allgemeine Theorie der biquadratischen Reste, wie die \textit{Comm. secunda} sie enthält. Wahrscheinlicher aber dünkt uns, dass \textsc{Gauss} jetzt aus dieser Theorie die Hilfsmittel gewonnen hat, um die (von dem Satze über den biquadratischen Charakter der Zahl 2 verschiedenen) Haupttheoreme zu beweisen, von denen er im oben erwähnten Briefe an \textsc{Dirichlet} schreibt, dass er ihre Beweise gerade etwa in dieser Zeit gefunden habe. Nach seiner Anzeige der \textit{Comm. secunda} vom 23. April 1831 (siehe W. II, S. 173) waren ihm freilich diese Beweise schon »seit 20 Jahren« bekannt, doch ist diese Angabe wohl nur als eine abgerundete zu betrachten. Vielleicht sind die Nummern 144, 145 des \textit{Tagebuchs} aber auch so zu verstehen, dass \textsc{Gauss} in der Kreisteilung, wenn nicht gar, wie \textsc{Eisenstein}\begingroup\renewcommand{\thefootnote}{**)}\footnote{\textsc{Crelles} Journal für Mathematik \textbf{30} (1846), S. 185.}\endgroup, in der Lemniskatenteilung die Hilfsmittel zu jenen Beweisen gefunden hat; für das letztere könnte die anschliessende letzte Tagebuchnotiz Nr. 146, die einzige, in welcher von den komplexen Zahlen \(a+bi\) die Rede ist, einigen Anhalt bieten.
